%========================================
% Section：[9] Taṅkhaṇika-paccavekkhaṇa-pāṭho
%========================================
% title-en: Reflection at the Time of Use
\section*{m9 | Taṅkhaṇika-paccavekkhaṇa-pāṭho}

\begin{chantsubsection}
  \chantnote
    {\{Handa mayaṃ taṅkhaṇika-paccavekkhaṇa-pāṭhaṃ bhaṇāmase\}}
  \chantnotetrans
    {\{Now let us recite the reflection at the time of using the requisites.\}}
  \chantgap

  % Cīvara (robe)
  \chantlineinline
    {Paṭisaṅkhā yoniso cīvaraṃ paṭisevāmi,}
    {Wisely reflecting, I use the robe,}
  \chantlineinline
    {yāvadeva sītassa paṭighātāya,}
    {only for warding off cold,}
  \chantlineinline
    {uṇhassa paṭighātāya,}
    {for warding off heat,}
  \chantlineinline
      {ḍaṃsa-makasa-vātātapa-siriṃsapa-samphassānaṃ paṭighātāya,}
    {for warding off the touch of flies, mosquitoes, wind, sun, and creeping things,}
  \chantlineinline
    {yāvadeva hirikopina-paṭicchādanatthaṃ.}
    {only for the purpose of covering the body.}

\chantgap

  % Piṇḍapāta (almsfood)
  \chantlineinline
    {Paṭisaṅkhā yoniso piṇḍapātaṃ paṭisevāmi,}
    {Wisely reflecting, I use almsfood:}
  \chantlineinline
    {neva davāya,}
    {not for fun,}
  \chantlineinline
    {na madāya,}
    {not for intoxication,}
  \chantlineinline
    {na maṇḍanāya,}
    {not for fattening,}
  \chantlineinline
    {na vibhūsanāya,}
    {not for beautification,}
  \chantlineinline
    {yāvadeva imassa kāyassa ṭhitiyā, yāpanāya,}
    {but only for the maintenance and nourishment of this body,}
  \chantlineinline
    {vihiṃsūparatiyā, brahmacariyānuggahāya,}
    {for keeping it unharmed and for supporting the holy life, thinking:}
  \chantlineinline
    {iti purāṇañca vedanaṃ paṭihaṅkhāmi,}
    {'thus I shall allay old feeling,}
  \chantlineinline
    {navañca vedanaṃ na uppādessāmi,}
    {and not arouse new feeling;}
  \chantlineinline
    {yātrā ca me bhavissati,}
    {there will be for me support for the journey of life,}
  \chantlineinline
    {anavajjatā ca phāsuvihāro cā'ti.}
    {and blamelessness and comfortable abiding.}

\end{chantsubsection}
\chantgap
\begin{chantsubsection}

  % Senāsana (lodging)
  \chantlineinline
    {Paṭisaṅkhā yoniso senāsanaṃ paṭisevāmi,}
    {Wisely reflecting, I use a lodging,}
  \chantlineinline
    {yāvadeva sītassa paṭighātāya,}
    {only for warding off cold,}
  \chantlineinline
    {uṇhassa paṭighātāya,}
    {for warding off heat,}
  \chantlineinline
    {ḍaṃsa-makasa-vātātapa-siriṃsapa-samphassānaṃ paṭighātāya,}
    {for warding off the touch of flies, mosquitoes, wind, sun, and creeping things,}
  \chantlineinline
    {yāvadeva utuparissaya vinodanaṃ paṭisallānārāmatthaṃ.}
    {only for dispelling the hardships of climate and for the enjoyment of seclusion.}

\chantgap

  % Bhesajja (medicinal requisites)
  \chantlineinline
    {Paṭisaṅkhā yoniso gilāna-paccaya-bhesajja-parikkhāraṃ paṭisevāmi,}
    {Wisely reflecting, I use medicinal requisites for the sick,}
  \chantlineinline
    {yāvadeva uppannānaṃ veyyābādhikānaṃ vedanānaṃ paṭighātāya,}
    {only for warding off painful feelings that have arisen,}
  \chantlineinline
    {abyāpajjha-paramatāyā'ti.}
    {and for the attainment of freedom from affliction.}
\end{chantsubsection}
